\subsection{可解李代数的嘉当子代数的共轭性}

设 g 是一个可解李代数。用 \(\mathcal{C}^{\infty}(\mathfrak{g})\) 表示 g 的下降中心列诸项的交（第一章，§1，第 5 节）。它是 g 的一个特征理想，并且是使 g/m 幂零的 g 的最小理想 m。由于 \(\mathcal{C}^{\infty}(\mathfrak{g}) \subset [\mathfrak{g}, \mathfrak{g}]\)，\(\mathcal{C}^{\infty}(\mathfrak{g})\) 是 g 的一个幂零理想（第一章，§5，第 3 节，定理 1 的推论 5）。由第 1 节的命题 1，诸 \(e^{ad x}\)（\(x \in \mathcal{C}^{\infty}(\mathfrak{g})\)）所成的集是 \(\operatorname{Aut}(\mathfrak{g})\) 的一个子群，包含在特殊自同构群内（第一章，§6，第 8 节，定义 6）。

定理 3. 设 \(\mathfrak{g}\) 是可解李代数，\(\mathfrak{h},\mathfrak{h}'\) 是 \(\mathfrak{g}\) 的嘉当子代数。则存在 \(x\in \mathcal{C}^{\infty}(\mathfrak{g})\) 使得 \(e^{\mathrm{ad}x}\mathfrak{h} = \mathfrak{h}'\)。

我们对 \(\dim \mathfrak{g}\) 作归纳，\(\mathfrak{g} = 0\) 的情形是平凡的。设 \(\mathfrak{n}\) 是 \(\mathfrak{g}\) 的一个极小非零交换理想。设 \(\varphi : \mathfrak{g} \to \mathfrak{g} / \mathfrak{n}\) 为典则态射。则有 \(\varphi(\mathcal{C}^{\infty} \mathfrak{g}) = \mathcal{C}^{\infty}(\mathfrak{g} / \mathfrak{n})\)（第一章，§1，第 5 节，命题 4）。由于 \(\varphi(\mathfrak{h})\) 与 \(\varphi(\mathfrak{h}')\) 都是 \(\mathfrak{g} / \mathfrak{n}\) 的嘉当子代数（§2，第 1 节，命题 4 的推论 2），由归纳假设，存在 \(x \in \mathcal{C}^{\infty}(\mathfrak{g})\) 使得 \(e^{\mathrm{ad} \varphi(x)}\varphi(\mathfrak{h}) = \varphi(\mathfrak{h}')\)。把 \(\mathfrak{h}\) 换成 \(e^{\mathrm{ad} x}\mathfrak{h}\)，即可假设 \(\varphi(\mathfrak{h}) = \varphi(\mathfrak{h}')\)，换言之，

\[
\mathfrak {h} + \mathfrak {n} = \mathfrak {h} ^ {\prime} + \mathfrak {n}.
\]

于是 \(\mathfrak{h}\) 与 \(\mathfrak{h}'\) 是 \(\mathfrak{h} + \mathfrak{n}\) 的嘉当子代数。若 \(\mathfrak{h} + \mathfrak{n} \neq \mathfrak{g}\)，则待证的断言由归纳假设得出。以下设 \(\mathfrak{h} + \mathfrak{n} = \mathfrak{h}' + \mathfrak{n} = \mathfrak{g}\)。

由 \(\mathfrak{n}\) 的极小性，有 \([\mathfrak{g},\mathfrak{n}] = \{0\}\) 或 \([\mathfrak{g},\mathfrak{n}] = \mathfrak{n}\)。若 \([\mathfrak{g},\mathfrak{n}] = \{0\}\)，则 \(\mathfrak{n} \subset \mathfrak{h}\) 且 \(\mathfrak{n} \subset \mathfrak{h}'\)（§2，第 1 节，命题 5），于是 \(\mathfrak{h} = \mathfrak{h} + \mathfrak{n} = \mathfrak{h}' + \mathfrak{n} = \mathfrak{h}'\)。剩下考虑 \([\mathfrak{g},\mathfrak{n}] = \mathfrak{n}\) 的情形，此时 \(\mathfrak{n} \subset \mathcal{C}^{\infty}(\mathfrak{g})\)。理想 \(\mathfrak{n}\) 是一个单 \(\mathfrak{g}\)-模；由于 \(\mathfrak{g} = \mathfrak{h} + \mathfrak{n}\)，且 \([\mathfrak{n},\mathfrak{n}] = \{0\}\)，可知 \(\mathfrak{n}\) 是一个单 \(\mathfrak{h}\)-模。若 \(\mathfrak{h} \cap \mathfrak{n} \neq \{0\}\)，则 \(\mathfrak{n} \subset \mathfrak{h}\)，于是 \(\mathfrak{g} = \mathfrak{h}\) 且 \(\mathfrak{h}' = \mathfrak{h}\)。现在设 \(\mathfrak{h} \cap \mathfrak{n} = \{0\}\)。则 \(\mathfrak{g} = \mathfrak{h} \oplus \mathfrak{n}\)，从而 \(\mathfrak{g} = \mathfrak{h}' \oplus \mathfrak{n}\)，因为 \(\mathfrak{h}\) 与 \(\mathfrak{h}'\) 有相同的维数。

对所有 \(x \in h\)，设 \(f(x)\) 是 n 中使 \(x - f(x) \in \mathfrak{h}'\) 的唯一元素；若 \(x, y \in h\)，

\[
[ x, y ] - [ x, f (y) ] - [ f (x), y ] = [ x - f (x), y - f (y) ] \in \mathfrak {h} ^ {\prime},
\]

于是 \(f([x,y]) = [x,f(y)] + [f(x),y]\)。由 §1 第 3 节命题 9 的推论，存在 \(a \in \mathfrak{n}\) 使得 \(f(x) = [x,a]\) 对所有 \(x \in \mathfrak{h}\) 成立。我们有 \((\operatorname{ad}a)^2(\mathfrak{g}) \subset (\operatorname{ad}a)(\mathfrak{n}) = 0\)，故对所有 \(x \in \mathfrak{h}\)，

\[
e ^ {\operatorname{ad} a} x = x + [ a, x ] = x - f (x).
\]

这样就有 \(e^{\mathrm{ad}a}(\mathfrak{h}) = \mathfrak{h}'\)。由于 \(a \in \mathcal{C}^{\infty}(\mathfrak{g})\)，证明完毕。

引理 3. 设 g 是李代数，r 是其根基，\(\varphi\) 是 g 到 g/r 的典则同态，v 是 g/r 的一个初等自同构。则存在 g 的初等自同构 u 使得 \(\varphi \circ u = v \circ \varphi\)。

我们可以假设 \(v\) 形如 \(e^{\mathrm{ad}b}\)，其中 \(b \in \mathfrak{g} / \mathfrak{r}\) 且 \(\mathrm{ad}b\) 幂零。设 \(\mathfrak{s}\) 是 \(\mathfrak{g}\) 的一个 Levi 子代数（第一章，§6，第 8 节，定义 7），并设 \(a\) 是 \(\mathfrak{s}\) 中使 \(\varphi(a) = b\) 的元素。由于 \(\mathrm{ad}_{\mathfrak{s}}a\) 幂零，\(\mathrm{ad}_{\mathfrak{g}}a\) 幂零（第一章，§6，第 3 节，命题 3 的推论），且 \(u = e^{\mathrm{ad}_{\mathfrak{g}}a}\) 是 \(\mathfrak{g}\) 的一个初等自同构，满足 \(\varphi \circ u = v \circ \varphi\)。

命题 5. 设 g 是李代数，r 是其根基，h 与 h' 是 g 的嘉当子代数，\(\varphi\) 是 g 到 g/r 的典则同态。下列条件等价：

\begin{enumerate}
\def\labelenumi{(\roman{enumi})}
\item
  \(\mathfrak{h}\) 与 \(\mathfrak{h}'\) 经 \(\mathfrak{g}\) 的一个初等自同构共轭；
\item
  \(\varphi (\mathfrak{h})\) 与 \(\varphi (\mathfrak{h}^{\prime})\) 经 \(\mathfrak{g} / \mathfrak{r}\) 的一个初等自同构共轭。
\item
  \(\Longrightarrow\) (ii)：这是显然的。
\item
  \(\Longrightarrow\) (i)：假设条件 (ii) 成立，证明 (i)。由引理 3，可化为 \(\varphi(\mathfrak{h}) = \varphi(\mathfrak{h}')\) 的情形。置 \(\mathfrak{k} = \mathfrak{h} + \mathfrak{r} = \mathfrak{h}' + \mathfrak{r}\)，它是 \(\mathfrak{g}\) 的一个可解子代数。于是 \(\mathfrak{h}\) 与 \(\mathfrak{h}'\) 是 \(\mathfrak{k}\) 的嘉当子代数，故存在 \(x \in \mathcal{C}^{\infty}(\mathfrak{k})\) 使得 \(e^{\mathrm{ad}_{\mathfrak{k}}x}\mathfrak{h} = \mathfrak{h}'\)（定理 3）。由于 \(\mathfrak{k}/\mathfrak{r}\) 幂零，\(\mathcal{C}^{\infty}(\mathfrak{k}) \subset \mathfrak{r}\)；另一方面，\(\mathcal{C}^{\infty}(\mathfrak{k}) \subset [\mathfrak{k},\mathfrak{k}] \subset [\mathfrak{g},\mathfrak{g}]\)，故 \(x \in \mathfrak{r} \cap [\mathfrak{g},\mathfrak{g}]\)；由第一章 §5 第 3 节定理 1，\(\mathrm{ad}_{\mathfrak{g}}x\) 幂零，从而 \(e^{\mathrm{ad}_{\mathfrak{g}}x}\) 是 \(\mathfrak{g}\) 的一个把 \(\mathfrak{h}\) 变换为 \(\mathfrak{h}'\) 的初等自同构。
\end{enumerate}

